\documentclass[12pt]{article}
\usepackage[margin=1in]{geometry}
\usepackage{fancyhdr}
\pagestyle{fancy}
\fancyhf{}
\fancyhead[L]{PD2MD Technical Manual}
\fancyhead[R]{Version 1.0}
\fancyfoot[C]{Page \thepage\ of 3}
\fancyfoot[R]{Confidential}
\title{Document with Headers and Footers}
\author{Header/Footer Test}
\begin{document}
\maketitle
\thispagestyle{fancy}

\section{Page 1 Content}
This document has running headers (``PD2MD Technical Manual'' on the left,
``Version 1.0'' on the right) and footers (page number centered, ``Confidential''
on the right). The converter should detect these repeating elements across pages
and exclude them from the body content in the Markdown output.

The header and footer text should not appear as paragraphs in the converted
Markdown. Only the body content within each section should be preserved.

\newpage
\section{Page 2 Content}
This is the second page of the document. The same headers and footers appear
here, demonstrating the cross-page repetition pattern that the converter uses
to identify running headers and footers.

Content on this page should flow naturally after the content from page 1 in the
Markdown output, without any header or footer text interrupting the flow.

\newpage
\section{Page 3 Content}
The third and final page continues with the same layout. By having three pages,
we provide enough data for the cross-page header/footer detection algorithm to
work reliably. The algorithm requires at least two pages to identify patterns.

All body content from all three pages should appear in the correct order in the
final Markdown output.
\end{document}